\honest{How An Algorithm Feels From Inside}{Eliezer Yudkowsky}{2008}

Two people argue about whether a tree that falls unheard makes a sound. One means vibrations in the air, the other means a heard experience, so the argument is about a word. I accept this standard view, and ask a better question: why do people argue about it at all?

I bring back a sorting task from an earlier post, ``Disguised Queries.'' Blue eggs, called ``bleggs,'' usually contain vanadium and glow in the dark; red cubes, called ``rubes,'' contain palladium. A blue egg with palladium in it goes in the rube bin but glows like a blegg. So ``Is it a blegg?'' depends on which practical question you mean, and once you know every property, no question is left.

Why, then, do people keep asking whether it is ``really'' a blegg? I show two network designs. In the first, every unit stands for an observable property. In the second, the properties all connect to a central unit, which can stay undecided after everything is known. The second is ``a far better candidate for being something vaguely like how the human brain works,'' because ``It's fast, cheap, scalable.'' \nb{Those are reasons an engineer might choose it. They are not evidence that brains use it.} I cite no study of how brains or people sort things into categories.

That central unit, I say, is what the leftover question feels like from the inside. Did the tree make a sound? Is Pluto a planet?

We see our intuitions as the world itself, I say. When you look at a green cup, you are really seeing ``a picture reconstructed in your visual cortex,'' though you think you see the cup. \nb{Philosophers have argued about exactly this for centuries; some hold that you see the cup, and that the picture is how you see it.} I settle the question in a dash.

Even calling the dispute over Pluto a matter of definitions, I say, is thinking like the second network. A mind built like the first would feel that no question was left.

I end by suggesting that people ``cling to their intuitions'' because they cannot see them as the way their own minds work, so that anything you say about how their thinking goes wrong gets ``discarded as obviously wrong.'' The words ``cling to their intuitions'' link to my post ``Allais Malaise,'' answering readers who had disagreed with me about a puzzle in decision theory. I give no way to tell a person clinging to an intuition from a person with a good objection.
